\section{Experimental Design}

The study uses a three-tier design that evaluates lightweight, midclass, and highclass local models under the same highway benchmark, then applies a post-hoc validity screen before making any comparative claim. This separates two questions that local LLM papers often blur together: how well a model drives when a run is valid, and how often the local runtime policy allows that run to remain valid at all. Figure~\ref{fig:benchmark-test-specification} summarizes the benchmark test specification itself: a fixed \texttt{lampilot\_highway\_v1} suite running on \metric{highway-fast-v0}, eight task categories, five cases per category, and a validity-aware reporting path that keeps invalid runs visible. Table~\ref{tab:validity-summary} and Figure~\ref{fig:cross-tier-validity} then provide the measured counts that instantiate this design. In the current bundle, lightweight has seven valid rows and one exact eligible base-versus-tuned pair, midclass has only three valid rows and no exact eligible pairs, and highclass has no valid rows at all. Those counts already determine what kind of inference each tier can support.

\begin{figure}[t]
\centering
\includegraphics[width=\linewidth]{benchmark_test_specification.png}
\caption{Benchmark-test specification for the highway study. The figure summarizes the fixed \texttt{lampilot\_highway\_v1} suite, its task taxonomy, and the validity-aware evaluation path used for all model tiers. It is a benchmark-design figure rather than a performance summary.}
\label{fig:benchmark-test-specification}
\end{figure}

Within the valid subset, our headline task metric is \metric{driving\_score\_v2}, reported alongside task completion, crash rate, and mean decision latency. We do not collapse these into a single global score, and we do not rank invalid runs by latency alone. Instead, the design is deliberately layered. First, we ask whether a tier produces enough valid rows to support ranking at all. Second, for tiers that do, we use valid-only leaderboard-style comparisons to characterize the operating region that the benchmark can still compare fairly. Third, we ask whether any exact base-versus-fine-tuned pairs survive that screen and therefore support a causal fine-tuning interpretation. This staged procedure makes the logic of the paper explicit: screening comes before ranking, and ranking comes before adaptation claims.

This is why Table~\ref{tab:lightweight-leaderboard} and Figure~\ref{fig:cross-tier-pareto} function as main-paper evidence, whereas the midclass and highclass summaries are interpreted as screening and failure-analysis results rather than as symmetric benchmark tables. The lightweight assets describe the only tier in which multiple valid local models can still be ranked and interpreted together. The midclass table records which families retain at least one valid representative while simultaneously documenting why exact family-level attribution remains blocked. The highclass summary has a different purpose again: it demonstrates that the current local runtime regime can exhaust the benchmark's interpretability before a family-level comparison ever begins.

The resulting design is asymmetric by construction. Lightweight is the only tier that supports both comparative ranking and an exact pairwise fine-tuning test, with the Phi pair providing the paper's direct adaptation evidence in Table~\ref{tab:lightweight-pair-delta}. Midclass still offers partial visibility into valid local behavior, but not enough matched coverage for clean pairwise attribution. Highclass serves a different scientific role: it shows where the current local runtime regime causes the benchmark to stop producing interpretable rows altogether. We now turn to that screening structure and the small set of comparisons it actually licenses.
